\documentclass[11pt]{article}

\usepackage[T1]{fontenc}
\usepackage[utf8]{inputenc}
\usepackage{amsmath,amssymb,amsthm}
\usepackage{booktabs}
\usepackage{geometry}
\usepackage{hyperref}
\usepackage{xcolor}
\usepackage{microtype}

\geometry{margin=1in}
\hypersetup{
  colorlinks=true,
  linkcolor=blue!55!black,
  citecolor=blue!55!black,
  urlcolor=blue!55!black
}

\newtheorem{theorem}{Theorem}
\newtheorem{proposition}[theorem]{Proposition}
\newtheorem{lemma}[theorem]{Lemma}
\newtheorem{corollary}[theorem]{Corollary}
\theoremstyle{definition}
\newtheorem{definition}[theorem]{Definition}
\theoremstyle{remark}
\newtheorem{remark}[theorem]{Remark}

\newcommand{\F}{\mathbb F}
\newcommand{\Dfive}{D_5}
\newcommand{\symdiff}{\mathbin{\triangle}}
\newcommand{\kap}{\kappa}
\newcommand{\wt}{\operatorname{wt}}

\title{Coloured Surfaces and Kempe Surgery for \(D_5\)-Flows:\\
Structural Lemmas, a Reconfiguration Trap, and Finite Evidence}
\author{Working research note for independent human verification}
\date{July 28, 2026}

\begin{document}
\maketitle

\begin{abstract}
The standard five-cycle-double-cover conjecture asks whether every finite
bridgeless graph has at most five Eulerian subgraphs covering each edge
exactly twice.  It remains unresolved.  This note records a structural
model and exact computations arising from an attempted proof.

For a loopless cubic graph, a five-cycle double cover is equivalently a
conservative labeling of every edge by a two-subset of five coordinates,
called a \(D_5\)-flow.  We show directly that such a flow canonically
determines a closed triangulated surface whose faces are dual to the graph
vertices and whose primal vertices are properly coloured with five
colours.  The coordinate circuits are precisely the primal vertex links,
and
\[
 \chi(S(q))=\sum_{i=0}^{4}\kap(C_i)-\frac{|V(G)|}{2}.
\]
The usual componentwise coordinate transposition is a reversible
triangle-band recolouring and regluing surgery.  We give an exact
permutation formula for its change in Euler characteristic and a cube
example in which one switch changes a torus into a sphere.

We also prove a four-coordinate locking lemma and give a separately
checkable 14-vertex contraction whose complete generalized-switch orbit
has twelve normalized states, none satisfying the boundary condition
needed to undo the contraction.  This refutes a natural reconfiguration
lemma, not the five-cycle-double-cover conjecture.  A complementary Tait
quotient proves that the analogous fully circuit-trapped branch is
root-universal on cubic graphs, explaining why the degree-four trap does
not lift unchanged.  We also construct a planar cubic capped-ladder
family proving that no constant number of root-component Kempe switches
can always rescue two prescribed edges; this concerns a restricted
strategy and does not refute rooted Kempe universality or FiveCDC.
Finally, exhaustive computations show that every
terminal equal-Euler-characteristic Kempe plateau is root-universal for
all 81 biconnected simple cubic graphs on 12 vertices.  A weaker
global-maximum test succeeds for all 480 such graphs on 14 vertices.
These are finite results only.  The universal terminal-plateau statement
is left explicitly as a conjecture.
\end{abstract}

\section{Status, conventions, and contribution boundary}

A \(k\)-cycle double cover, in the convention used here, is a list of at
most \(k\) Eulerian edge-subsets such that every edge occurs in exactly two
members.  Empty members may be appended, so ``at most five'' and ``five,
allowing empty members'' are equivalent.  This agrees with the convention
in Oum's recent exposition: the general cycle-double-cover theorem gives
eight Eulerian subgraphs, while the reduction from eight to five remains
Conjecture~18~\cite{Oum2026}.  Recent work on five-cycle double covers also
continues to state the five-subgraph assertion as a conjecture and proves
special cases~\cite{LiuEtAl2023,LiuEtAl2025}.

The orientable five-cycle-double-cover conjecture adds directed Eulerian
orientations so that each edge is used once in each direction.  It is a
different, stronger statement; no orientation condition is imposed in
this note.

The embedding aspect is established background, not a novelty claim.
Catlin explicitly studied embedded graphs, facial colourings, and double
cycle covers~\cite{Catlin1990}; see also Zhang's monograph
\cite{Zhang2012} and Mohar and Thomassen's text
\cite{MoharThomassen2001}.  Ghanbari and \v{S}\'amal formulate the cubic
CDC in terms of embeddings without singular edges and study both
approximation and twist operations~\cite{GhanbariSamal2025,
GhanbariSamal2026}.  We have not completed a publication-grade prior-art
search for the exact \(D_5\) five-colour specialization, the
component-switch Euler formula below, or the particular finite trap.
Accordingly, this draft makes no categorical novelty claim.  Its present
value is that the statements, proofs, certificates, and limitations are
exposed in a form that can be checked independently.

\textbf{Literature-search TODO.}
Before circulation as a preprint, conduct a targeted search for prior
work on unbounded Kempe-switch distance in two-edge-cut chains, especially
prefix/suffix word obstructions for cycle-double-cover or nowhere-zero
flow reconfiguration.  The capped-ladder theorem below has not yet been
compared adequately against that literature.

\section{\texorpdfstring{\(D_5\)-flows}{D5-flows}}

All graphs in the structural sections are finite and loopless.  Parallel
edges are permitted unless simplicity is stated.  Edge objects, including
parallel ones, are distinct.

Put
\[
 \Dfive=\binom{\{0,1,2,3,4\}}2\subset\F_2^5.
\]
We identify a pair \(ij\) with its weight-two incidence vector.  A
\emph{\(D_5\)-flow} on a graph \(G\) is a map
\[
 q:E(G)\longrightarrow D_5
\]
such that
\[
 \bigoplus_{e\ni v}q(e)=0                                    \tag{1}
\]
at every vertex \(v\).  Since the graphs considered here are loopless,
each incident edge appears once in (1).

\begin{proposition}[cover--flow equivalence]\label{prop:cover-flow}
The \(D_5\)-flows on \(G\) are in bijection with ordered five-tuples
\((C_0,\ldots,C_4)\) of Eulerian edge-subsets in which every edge occurs
exactly twice.
\end{proposition}

\begin{proof}
Given \(q\), put
\[
 C_i=\{e\in E(G):i\in q(e)\}.
\]
The \(i\)-th coordinate of (1) says that \(C_i\) has even degree at every
vertex.  Every label has weight two, so every edge belongs to exactly two
of the \(C_i\).

Conversely, given the five Eulerian subsets, label an edge by the two
indices of the members containing it.  The label has weight two, and the
five even-degree conditions are exactly (1).  The constructions are
inverse.
\end{proof}

The following elementary local observation drives the surface model.

\begin{lemma}[local coordinate triangle]\label{lem:local-triangle}
Let \(a,b,c\in D_5\) and \(a+b+c=0\).  Then there are three distinct
coordinates \(i,j,k\) for which
\[
 \{a,b,c\}=\{ij,ik,jk\}.
\]
\end{lemma}

\begin{proof}
We have \(c=a+b\).  Since
\[
 \wt(a+b)=4-2|a\cap b|=2,
\]
the two pairs \(a,b\) meet in exactly one coordinate.  Their union has
three coordinates, and \(c=a\symdiff b\) is the third side of that
coordinate triangle.
\end{proof}

\section{The coloured-surface dictionary}

We first give a construction that does not assume an embedding in
advance.

\begin{theorem}[canonical coloured surface]\label{thm:surface}
Let \(G\) be a connected loopless cubic graph with a \(D_5\)-flow \(q\).
Then \(q\) canonically determines, up to colour-preserving cellular
isomorphism, a connected closed triangulated surface \(S(q)\) with the
following properties.
\begin{enumerate}
 \item The triangles of \(S(q)\) are indexed by \(V(G)\), and the dual
 graph is \(G\).
 \item The primal vertices are properly coloured with coordinates
 \(0,\ldots,4\), in the sense that the three corners of every triangle
 have distinct colours.
 \item If \(C_i=\{e:i\in q(e)\}\), the components of \(C_i\) are in
 bijection with the primal vertices of colour \(i\).
\end{enumerate}
Consequently, if \(n=|V(G)|\), then
\[
 \chi(S(q))=\sum_{i=0}^4\kap(C_i)-\frac n2.                  \tag{2}
\]
\end{theorem}

\begin{proof}
For each \(v\in V(G)\), Lemma~\ref{lem:local-triangle} gives three
coordinate colours on the incident labels.  Take an abstract triangle
\(\Delta_v\) with those vertex colours, and assign the side opposite the
remaining colour to the incident graph edge whose label is the two
endpoint colours of that side.

For every graph edge \(e=uv\), glue the side of \(\Delta_u\) indexed by
\(e\) to the side of \(\Delta_v\) indexed by \(e\), using the unique map
that preserves its two endpoint colours.  Every triangle side is paired
exactly once.  The quotient is therefore a closed two-dimensional
pseudomanifold.

A quotient vertex can initially have more than one link component.
Replace it by one vertex for each link component.  Each resulting vertex
has a circular link, so the normalized quotient is a closed triangulated
surface.  Splitting vertices does not change the triangles, their side
pairings, or the dual graph.  The surface is connected because its dual
graph \(G\) is connected.

Fix a coordinate \(i\).  In a triangle having an \(i\)-coloured corner,
the two sides at that corner are exactly the two incident graph edges
whose labels contain \(i\).  Following the link of that corner through the
side pairings therefore follows a component of \(C_i\), and every such
component arises this way.  This proves (1)--(3).

The surface has \(n\) triangles and \(3n/2\) primal edges.  Its number of
vertices is \(\sum_i\kap(C_i)\).  Euler's formula now gives
\[
 \chi(S(q))=\sum_i\kap(C_i)-\frac{3n}{2}+n,
\]
which is (2).
\end{proof}

For a precise converse, call a face-colouring of a cellular embedding of
a cubic graph \emph{corner-proper} when the three face incidences at each
graph vertex have three distinct colours.  This definition rules out a
single face occupying two corners at one vertex.

\begin{proposition}[converse]\label{prop:surface-converse}
A corner-proper five-colouring of the faces of a cellular embedding of a
connected cubic graph gives a \(D_5\)-flow: label an edge by the colours of
its two incident face occurrences.  This construction and
Theorem~\ref{thm:surface} are inverse after normalizing quotient vertices
by link components.
\end{proposition}

\begin{proof}
At a cubic vertex, let the three incident face colours be \(i,j,k\).  The
three incident edge labels are \(ij,ik,jk\), whose xor is zero.  Thus the
labels form a \(D_5\)-flow.  Taking the dual triangle at every graph vertex
recovers the original side pairings and face-colour occurrences.
Normalization separates exactly the distinct link components represented
by distinct facial occurrences.
\end{proof}

\begin{remark}
Propositions~\ref{prop:cover-flow} and~\ref{prop:surface-converse} show
that the standard, unoriented five-cycle-double-cover problem for cubic
graphs can be read as a problem of finding a corner-proper
five-face-coloured cellular embedding.  This is not an orientability
statement: \(S(q)\) may be orientable or nonorientable.
\end{remark}

\section{Bichromatic boundary curves and Kempe surgery}

For \(i\ne j\), define
\[
 Y_{ij}=C_i\symdiff C_j
       =\{e:|q(e)\cap\{i,j\}|=1\}.                            \tag{3}
\]
It is an even subgraph and, in a cubic graph, a disjoint union of
circuits.

\begin{proposition}[boundary interpretation]\label{prop:boundary}
In \(S(q)\), let \(H_{ij}\) be the primal subgraph induced by the primal
vertices of colours \(i,j\).  After the standard dual isotopy,
\[
 \partial N(H_{ij})=Y_{ij},                                  \tag{4}
\]
where \(N(H_{ij})\) is a sufficiently small regular neighbourhood.
\end{proposition}

\begin{proof}
Check one primal triangle.  If it contains neither \(i\) nor \(j\), the
neighbourhood has no boundary arc there.  If it contains exactly one of
them, the boundary arc around that corner crosses the two triangle sides
whose labels contain exactly one of \(i,j\).  If it contains both, the
boundary arc around their common primal edge crosses the other two sides,
again exactly those whose labels contain one of \(i,j\).  These local arcs
join across paired triangle sides and give (4).
\end{proof}

Let \(K\) be a circuit component of \(Y_{ij}\).  A \emph{component Kempe
switch} transposes \(i,j\) in the label of every edge of \(K\).

\begin{lemma}[validity of a component switch]\label{lem:valid-switch}
A component Kempe switch sends a \(D_5\)-flow to a \(D_5\)-flow and is an
involution.
\end{lemma}

\begin{proof}
For an active label, transposing \(i,j\) adds the vector \(ij\).  At each
graph vertex the circuit \(K\) has degree zero or two, so the total added
vector in the local xor equation is zero.  The switched labels still have
weight two.  Applying the same transposition a second time restores the
original state.
\end{proof}

On the coloured surface, this move has a concrete description.  Recolour
the triangles traversed by \(K\) by the transposition \(i\leftrightarrow
j\).  Active side gluings still match because both incident triangles
are recoloured.  At an inactive \(ij\)-side incident with exactly one
recoloured triangle, compose the side gluing with the endpoint flip.  The
result is exactly the coloured quotient for the switched flow.  Thus the
move is a reversible triangle-band recolouring and regluing surgery.
It need not act on one fixed topological surface.

\subsection{An exact Euler-change formula}

Fix \(i,j\), and abbreviate \(A=C_i\), \(B=C_j\).  At each vertex of
\(A\cup B\), the nonzero membership patterns in \((A,B)\) are
\[
 10,\qquad 01,\qquad 11.
\]
Discard components of \(A\cup B\) that are pure circuits.  In every
remaining component suppress the degree-two paths.  The result is a cubic
multigraph \(T_{ij}\) whose three edge classes \(10,01,11\) are perfect
matchings.  Write \(m_{10},m_{01},m_{11}\) for the associated
fixed-point-free involutions on \(V(T_{ij})\).

For a permutation \(\pi\), let \(z(\pi)\) denote its number of cycles,
including fixed points.

\begin{theorem}[one-switch Euler formula]\label{thm:euler-change}
Let \(K\) be a component of \(Y_{ij}\) meeting \(T_{ij}\), and let \(W\)
be the vertices on the corresponding alternating circuit of
\(m_{10}\cup m_{01}\).  Define
\[
\begin{array}{c|cc}
 &W&V(T_{ij})\setminus W\\ \hline
m'_{10}&m_{01}&m_{10}\\
m'_{01}&m_{10}&m_{01},
\end{array}
\qquad m'_{11}=m_{11}.                                      \tag{5}
\]
If \(q'\) is obtained by switching on \(K\), then
\begin{equation}
\begin{split}
\chi(S(q'))-\chi(S(q))
 =\frac12\bigl[&
 z(m'_{10}m_{11})+z(m'_{01}m_{11})\\
 &-z(m_{10}m_{11})-z(m_{01}m_{11})\bigr].
\end{split}
\tag{6}
\end{equation}
If \(K\) is a discarded pure circuit, the Euler change is zero.
\end{theorem}

\begin{proof}
Only \(C_i=A\) and \(C_j=B\) can change, so by (2) the Euler change is
\[
 \Delta\kap(A)+\Delta\kap(B).                               \tag{7}
\]
The components of \(A\) that meet the cubic core correspond to the
alternating circuits of the two matchings \(m_{10},m_{11}\); similarly,
the core components of \(B\) correspond to the alternating circuits of
\(m_{01},m_{11}\).

If \(m_a,m_b\) are two perfect matchings, then on each alternating circuit
of their union the product \(m_am_b\) has two permutation cycles.
Therefore the number of alternating circuits is
\(\tfrac12z(m_am_b)\).

Switching \(i,j\) on \(K\) exchanges the \(10\)- and \(01\)-matching
edges at exactly the vertices of \(W\), giving (5).  Substitute the
before-and-after alternating-circuit counts into (7) to obtain (6).
Pure circuit components outside the cubic core are unchanged or move
whole from \(A\) to \(B\); their total contribution to (7) cancels.
\end{proof}

\subsection{A hand-checkable topology change}

Take the cube with edge order
\[
\begin{array}{c|rrrrrrrrrrrr}
e&0&1&2&3&4&5&6&7&8&9&10&11\\ \hline
uv&01&03&04&12&17&23&26&35&45&47&56&67
\end{array}
\]
and labels
\[
 01,02,12,02,12,01,12,12,01,02,02,01.                      \tag{8}
\]
The reader can check the xor equation at each of the eight vertices
directly from the table.  The coordinate component counts are
\[
 (\kap(C_0),\ldots,\kap(C_4))=(2,1,1,0,0),
\]
so (2) gives \(\chi=0\).  The edges \(\{0,1,3,5\}\) form one
\(Y_{12}\)-circuit.  Transposing \(1,2\) on those four edges changes the
component counts to
\[
 (2,2,2,0,0),
\]
so the new connected surface has \(\chi=2\).  A triangle-orientation
propagation check shows that both surfaces are orientable.  Hence this
single legal switch changes a torus into a sphere.  In particular, a
global proof cannot treat all Kempe moves as curve slides on one fixed
surface.

\section{A four-coordinate locking mechanism}

It is useful to distinguish component switches from a stronger move.
For a coordinate pair \(p=ij\), a \emph{generalized factor switch}
chooses any Eulerian edge-subset \(Z\subseteq Y_{ij}\) and adds the vector
\(ij\) to every label on \(Z\).  This remains a \(D_5\)-flow by the same
parity argument as Lemma~\ref{lem:valid-switch}.

\begin{lemma}[four-coordinate one-step lock]\label{lem:four-lock}
Suppose a \(D_5\)-flow uses exactly the four coordinates
\(S=\{0,1,2,3,4\}\setminus\{t\}\), and suppose every nonempty coordinate
class \(C_i\), \(i\in S\), is one circuit.  From this state, no generalized
factor switch can introduce \(t\) nontrivially modulo a global coordinate
permutation.
\end{lemma}

\begin{proof}
A switch with pair \(p\subseteq S\) does not introduce \(t\).  For a pair
\(p=\{i,t\}\), the absence of \(t\) gives
\[
 Y_{it}=C_i.
\]
A circuit has only two Eulerian edge-subsets, the empty set and the whole
circuit.  The empty switch does nothing.  Switching all of \(C_i\)
replaces \(i\) by \(t\) on every edge label containing \(i\), which is
exactly the global coordinate transposition \(i\leftrightarrow t\).
\end{proof}

The hypothesis of Lemma~\ref{lem:four-lock} need not be preserved by all
switches among the used coordinates.  Thus the lemma alone is a one-step
statement, not an orbit theorem.  The following finite example separately
certifies orbit closure.

\begin{proposition}[certified generalized-switch trap]\label{prop:trap}
There is a simple 3-connected cubic graph \(G\) on 14 vertices and an edge
\(e\) such that, after contracting \(e\), the resulting degree-four graph
has a generalized-factor-switch orbit with the following properties,
modulo the global \(S_5\) action:
\begin{enumerate}
 \item the orbit has exactly 12 states;
 \item every state uses four coordinates and each used coordinate class
 is one circuit; and
 \item every state has side xor of weight four and therefore cannot be
 extended over the restored edge.
\end{enumerate}
The same contracted graph also has a different orbit containing
extendible states.  Therefore this is not a counterexample to FiveCDC.
\end{proposition}

\begin{proof}[Finite certificate]
The parent graph has graph6 string
\[
 \texttt{M?AA@BORD\_CoEOAo?}.
\]
In lexicographic edge order, contract edge \(15=(3,12)\).  With the
contracted edges indexed by the inherited order, the prescribed sides at
the degree-four vertex are
\[
 U=\{5,11\},\qquad V=\{15,16\}.
\]
One state in the orbit is
\[
\begin{split}
(&01,02,03,13,03,01,23,02,03,03,02,23,23,13,12,\\
 &01,23,12,01,02).
\end{split}                                                  \tag{9}
\]
Its local xors vanish, while both side xors are
\[
 q(5)+q(11)=q(15)+q(16)=0123.                               \tag{10}
\]

The companion file \texttt{orbit-certificate.txt} lists all twelve
normalized rows and the complete normalized adjacency relation.  A
finite hand verification proceeds as follows.  For each row, check the
local xor equations and (10).  For each coordinate pair, construct
\(Y_{ij}\), delete one edge from every circuit to obtain a spanning
forest, and use the remaining edges as a binary cycle-space basis.
Switch every basis sum and normalize under the 120 coordinate
permutations.  The resulting distinct rows are exactly the neighbours
listed in the certificate.  The adjacency list is connected, so the
twelve rows form one complete orbit.  Each row also directly satisfies
the hypotheses of Lemma~\ref{lem:four-lock}.

The supplied independent Python checker performs exactly these finite
operations without a SAT solver, graph package, or canonical-labelling
library.  It also checks cubicity, bridgelessness, and 3-connectivity of
the parent.
\end{proof}

The proposition refutes the tempting assertion that every starting flow
on a contraction can be repaired by generalized factor switches.  It
does not refute the existential assertion that some flow on the
contraction extends; in fact, the other orbit contains 968 normalized
extendible states.

\subsection{Why the degree-four trap does not lift unchanged}

The contraction trap raises a natural concern: perhaps a cubic Kempe orbit
could remain locked in four coordinates with all four coordinate classes
connected and still fail rooted universality.  The next theorem excludes
exactly that branch.

Call a component-Kempe orbit on a cubic graph
\emph{four-coordinate circuit-trapped} when every state uses exactly four
coordinates and every used coordinate class is one circuit.

\begin{theorem}[cubic \(D_4\) trap theorem]\label{thm:d4-trap}
Every four-coordinate circuit-trapped component-Kempe orbit on a cubic
graph is root-universal.  More precisely, every pair of root edges is
rescued in the starting state or after one component switch between used
coordinates.
\end{theorem}

\begin{proof}
Fix a state \(q\), write \(S\) for its four used coordinates and \(t\) for
the missing coordinate, and let \(r,s\) be distinct root edges.  For
clarity, globally rename \(S\) as \(\{0,1,2,3\}\).

Identify complementary labels in the three classes
\[
 \{01,23\},\qquad \{02,13\},\qquad \{03,12\}.                \tag{11}
\]
These are three Tait colours.  At each cubic vertex, the local xor triple
is the triangle on three of the four coordinates, so its three labels lie
in distinct classes in (11).  Thus the quotient is a proper
three-edge-colouring.

If \(q(r)\cap q(s)\) contains a coordinate \(i\), then both roots lie in
\[
 C_i=Y_{it}.
\]
The circuit-trap hypothesis says that \(C_i\) is one circuit, so the roots
are already rescued.

Otherwise the two root labels are disjoint.  They are complementary
two-subsets of \(S\), hence have the same Tait colour \(A\).  Choose either
other Tait colour \(B\).  The union of classes \(A\cup B\) is a factor
\(Y_p\), where \(\{p,S\setminus p\}\) is the third class.  If \(r,s\) lie
on the same component of \(Y_p\), they are already rescued.

If they lie on different components, transpose the two coordinates of
\(p\) on the component containing \(r\).  In the Tait quotient this
interchanges the two active colours on that component: the colour of \(r\)
changes from \(A\) to \(B\), while the colour of \(s\) remains \(A\).
Labels from distinct complementary classes in (11) meet in exactly one
coordinate, say \(k\).  The literal switch did not introduce \(t\), so
both roots now lie in \(C'_k=Y'_{kt}\).  The switched state remains in the
assumed circuit-trapped orbit, and therefore \(C'_k\) is one circuit
containing both roots.
\end{proof}

The proof uses cubicity to obtain a proper Tait quotient and alternating
bichromatic factor circuits.  Neither fact survives at the degree-four
contraction vertex.  Theorem~\ref{thm:d4-trap} therefore explains why
Proposition~\ref{prop:trap} cannot be transplanted unchanged to a cubic
counterexample.  It closes only the fully circuit-trapped four-coordinate
branch; a general rooted-orbit obstruction could still use five
coordinates, a disconnected coordinate class, or transitions between
those regimes.

\section{Unbounded delay for root-component rescue}

Fix distinct root edges \(r,s\).  Call a state \emph{root-good} when some
component of some \(Y_{ij}\) contains both roots.  While the state is not
root-good, a \emph{root-component switch} means a component Kempe switch
whose switched circuit contains exactly one root.

For \(m\geq1\), define a capped ladder \(L_m\) as follows.  Its left cap
has vertices \(a,b,p,q\), its ladder has vertices \(u_i,v_i\) for
\(0\leq i\leq m\), and its right cap has vertices \(c,d,x,y\).  The edge
set is
\[
\begin{split}
 &ap,aq,bp,bq,pq,\quad au_0,bv_0,\\
 &u_iv_i\quad(0\leq i\leq m),\\
 &u_iu_{i+1},v_iv_{i+1}\quad(0\leq i<m),\\
 &u_mc,v_md,\quad cx,cy,dx,dy,xy.
\end{split}
\]
Thus the caps are copies of \(K_4-e\), and
\[
                  |V(L_m)|=2m+10,\qquad |E(L_m)|=3m+15.
\]
The graphs are simple planar cubic graphs.  They are bridgeless: cap
edges lie on cap triangles, ladder edges lie on squares, and the
attachment edges lie on cycles running through both rails and the caps.

Put
\[
                         A=01,\qquad B=02,\qquad C=12.
\]
There is a Tait-supported \(D_5\)-flow on \(L_m\) for which the labels on
the successive two-edge cuts from left to right form the length-\((m+2)\)
prefix
\[
                         A,C,B,A,C,B,\ldots.
\]
Explicitly, label \(au_0,bv_0\) by \(A\), label \(u_0v_0\) by \(B\),
and assign the rail-pair and rung labels successively so that every
ladder vertex sees \(A,B,C\).  On the left cap take
\[
 q(ap)=q(bq)=B,\quad q(aq)=q(bp)=C,\quad q(pq)=A,
\]
and use the unique compatible Tait colouring on the right cap, up to
exchanging its two inner vertices.  Every vertex then sees \(A,B,C\),
so the local xor equations hold.

\begin{theorem}[no constant root-component radius]
\label{thm:unbounded-root-radius}
Choose one root edge inside each cap of \(L_m\) in the state above.
Every root-component-switch sequence that reaches a root-good state has
length at least
\[
                         \left\lfloor\frac{m+2}{3}\right\rfloor.
\]
Consequently, no universal constant bounds the number of
root-component switches required for rescue.
\end{theorem}

\begin{proof}
Let \(D_0,\ldots,D_{m+1}\) be the nested two-edge cuts in the displayed
word.  For an arbitrary \(D_5\)-flow, the two edges of each \(D_j\) have
the same label: xor the vertex equations on either side of the cut;
internal edges cancel twice, leaving the xor of the two cut labels.
Write the resulting cut-label word as
\(\lambda_0\lambda_1\cdots\lambda_{m+1}\).

Let \(K\) be a factor component containing exactly one root.  Connectedness
and nestedness imply that the cuts met by \(K\) form a prefix when \(K\)
contains the left root and a suffix when it contains the right root.
Moreover, \(K\) is Eulerian, so it meets a two-edge cut in zero or two
edges.  A switch on \(K\) therefore applies one coordinate transposition
to a prefix or suffix of the cut-label word.

After \(t\) switches, their at most \(t\) boundaries divide the initial
word into at most \(t+1\) intervals.  On each interval the final labels
are obtained from the initial labels by one fixed coordinate permutation,
namely the ordered composition of the switches covering that interval.

If the final state is root-good through \(Y_h\), its component containing
both roots crosses every \(D_j\).  Hence
\[
                              h\cdot\lambda_j=1
                 \qquad(0\leq j\leq m+1),
\]
with dot product over \(\F_2\).

Partition the initial word into
\(\lfloor(m+2)/3\rfloor\) disjoint consecutive triples \(A,C,B\).
No one triple can remain in a single interval.  Otherwise one coordinate
permutation \(\pi\) would act on all three entries, and
\[
                \pi(A)+\pi(C)+\pi(B)=\pi(A+C+B)=0.
\]
Dotting with \(h\) contradicts the preceding display, because the xor of
three ones is one, not zero.  Thus every disjoint triple contains a
switch boundary in one of its two internal gaps.  Those gaps are disjoint
between triples, so one boundary splits at most one triple.  Therefore
\[
                         t\geq\left\lfloor\frac{m+2}{3}\right\rfloor,
\]
as required.
\end{proof}

\begin{remark}[scope]
The proof deliberately relaxes the actual move system and uses only its
necessary prefix/suffix action.  It proves unbounded distance for the
restricted root-component strategy.  It does not produce a state that is
unrescuable by root-component switches, does not refute unrestricted
rooted Kempe universality, and does not prove or disprove FiveCDC.
Literal breadth-first searches find finite rescue distances
\(3,4,5,6,8,9\) for selected roots in \(L_1,\ldots,L_6\), respectively.
These computations are supplementary to, not inputs to, the proof.
\end{remark}

\section{Surface-Euler plateaus}

Let \(\mathcal K(G)\) be the state graph whose vertices are \(D_5\)-flows
modulo global coordinate permutations and whose edges are component
Kempe switches.  Put
\[
 \Phi(q)=\chi(S(q)).
\]
For fixed graph edges \(r,s\), call \(q\) \emph{\((r,s)\)-good} if some
\(Y_{ij}\)-component contains both roots.

\begin{definition}
A \(\Phi\)-plateau is a connected component after retaining only the
equal-\(\Phi\) edges of \(\mathcal K(G)\).  It is \emph{terminal} when no
state in it has a neighbour of larger \(\Phi\).  A set of states is
\emph{root-universal} when, for every unordered edge pair \(r,s\), it
contains an \((r,s)\)-good state.
\end{definition}

\begin{lemma}[connected-coordinate non-descent]\label{lem:connected-coordinates}
Suppose every \(C_i\) is one nonempty circuit.  Then every component Kempe
switch \(q\mapsto q'\) satisfies
\[
 \Phi(q')\geq\Phi(q).
\]
If \(q\) lies in a terminal plateau, every such first switch is
\(\Phi\)-neutral.
\end{lemma}

\begin{proof}
Only the two switched coordinate classes, say \(C_i,C_j\), can change.
Initially their total component count is two.  If both resulting classes
are nonempty, each has at least one component and their total remains at
least two.

Suppose instead that \(C'_i\) is empty.  Then the switched component
contains all of \(C_i\), and \(C_i\cap C_j=\varnothing\), since intersection
edges are not in \(Y_{ij}\).  The two original circuits are also
vertex-disjoint: two edge-disjoint even subgraphs of a cubic graph cannot
both have degree two at the same vertex.  The new \(C'_j\) therefore
contains the two circuits as distinct components.  Its component count is
at least two, while \(C'_i\) contributes zero.  The case \(C'_j=\varnothing\)
is symmetric.  Thus
\[
 \kap(C'_i)+\kap(C'_j)\geq
 \kap(C_i)+\kap(C_j)=2,
\]
and (2) proves non-descent.  Terminality excludes a strict increase.
\end{proof}

\begin{lemma}\label{lem:reach-terminal}
Every state has a \(\Phi\)-nondecreasing Kempe path to a terminal plateau.
\end{lemma}

\begin{proof}
Move freely inside the current equal-\(\Phi\) plateau.  If it is not
terminal, reach a state having a neighbour of strictly larger \(\Phi\)
and take that edge.  The state graph is finite, so \(\Phi\) can increase
only finitely often.  The process ends at a terminal plateau.
\end{proof}

The computationally suggested statement is deliberately isolated.

\begin{quote}
\textbf{Terminal-plateau conjecture.}
For every connected bridgeless cubic graph with a \(D_5\)-flow, every
terminal \(\Phi\)-plateau is root-universal.
\end{quote}

This statement is unproved.  It is stronger than the finite observations
below and is not asserted as a theorem.

\subsection{Exact finite results}

The state enumeration quotients only by the global \(S_5\) action.  Graphs
were generated canonically by \texttt{geng -Cq -d3 -D3 n}.  The complete
order-12 terminal-plateau audit gives:
\[
\begin{array}{lr}
\toprule
\text{biconnected simple cubic graphs} &81\\
\text{normalized \(D_5\)-flows} &25{,}960\\
\text{Kempe orbits} &2{,}650\\
\text{terminal \(\Phi\)-plateaus} &2{,}659\\
\text{non-root-universal terminal plateaus} &0\\
\bottomrule
\end{array}                                                  \tag{12}
\]
The terminal values range from \(-3\) to \(2\), and plateau sizes range
from \(1\) to \(432\).  This is a complete finite census, not a proof of
the terminal-plateau conjecture.

An additional support diagnostic found 420 terminal plateaus in which
every state uses all five coordinates.  Thus the order-12 positive result
cannot be explained by always reaching a plateau with an unused
coordinate and reducing to a four-coordinate argument.  This support
diagnostic is present in the producer report but is not among the fields
recomputed by the focused independent verifier described below.

At order 14, a complete but weaker test was run.  In each Kempe orbit,
retain the states having the orbit's global maximum \(\Phi\), and split
them into connected equal-\(\Phi\) components.  Every such component was
root-universal:
\[
\begin{array}{lr}
\toprule
\text{biconnected simple cubic graphs} &480\\
\text{normalized \(D_5\)-flows} &537{,}418\\
\text{Kempe orbits} &33{,}102\\
\text{maximum number of max-\(\Phi\) components in one orbit} &12\\
\text{non-root-universal max-\(\Phi\) components} &0\\
\bottomrule
\end{array}                                                  \tag{13}
\]
Equation (13) does not test terminal plateaus below the global orbit
maximum and must not be reported as a complete order-14 test of the
terminal-plateau conjecture.

As a separate control, the union of each entire Kempe orbit was tested for
root-universality through order 14.  The same 33,102 order-14 orbits gave
6,951,420 orbit/root-pair tests and zero failures.  Again, this is finite
evidence only.

\subsection{Reproducibility and independent checks}

The complete reports, producers, and focused independent verifiers are in
the companion repository.  From its root:
\begin{verbatim}
python3 scratch/audit_d5_surface_euler_switch.py
python3 scratch/check_d5_contraction_circuit_orbit_counterexample.py
python3 scratch/check_d4_cubic_trap_root_universality.py
python3 scratch/check_d5_capped_ladder_delay_family.py

python3 scratch/verify_d5_surface_chi_plateaus_order12.py
python3 scratch/verify_d5_root_euler_potential_reports.py
python3 scratch/verify_d5_root_kempe_order14_report.py
\end{verbatim}
The order-12 plateau verifier regenerates all graph identities and report
arithmetic, then independently recomputes five deterministic
boundary/extremal rows with a separately structured traversal.  The
order-14 maximum-\(\Phi\) verifier regenerates all 480 graph identities
and report arithmetic, then independently recomputes ten selected rows
across orders 12 and 14.  These focused replays are independent semantic
checks, not second full enumerations; that limitation is part of the
claim.

\section{What has and has not been established}

The human-checkable results in this note are:
\begin{enumerate}
 \item the \(D_5\)-flow/five-even-subgraph equivalence;
 \item the canonical properly five-coloured surface construction and
 Euler formula;
 \item the boundary interpretation of the bichromatic factors;
 \item the exact one-switch Euler-change formula; and
 \item the four-coordinate one-step locking lemma and the cubic \(D_4\)
 trap theorem; and
 \item the capped-ladder lower bound proving that root-component rescue
 has no universal constant radius.
\end{enumerate}
The cube topology change is a short explicit certificate.  The
12-state contraction trap and the plateau tables are finite
computer-assisted results with full instances, producers, and independent
focused checkers.

No proof or disproof of the standard FiveCDC has been obtained.  No
orientable FiveCDC result has been obtained.  The terminal-plateau
conjecture is a proof target suggested by the computations, not a theorem.
A responsible public submission would require independent checking by a
graph theorist, a broader literature comparison, and archival publication
of the exact code and reports.

\section*{AI-use disclosure}

This project used OpenAI Codex agents extensively under human direction.
The agents proposed the \(D_5\) and coloured-surface viewpoints, derived
candidate lemmas and proofs, wrote and ran the enumeration programs,
searched for counterexamples to intermediate claims, produced the finite
certificates, and drafted this manuscript.  The human operator selected
the research objective and directed the work.  At the date of this draft,
the manuscript has not received independent expert verification or peer
review.

AI systems are not listed as authors.  Any human who circulates or submits
this note must personally verify the mathematical claims, take
responsibility for the text and computations, confirm compliance with the
venue's AI policy, and replace the working-draft author line with the
responsible author list.  The disclosure is intentionally broad: it does
not characterize the AI as a mere copy editor.

\bibliographystyle{plain}
\bibliography{references}

\end{document}
